% Part: normal-modal-logic
% Chapter: completeness
% Section: lindenbaums-lemma

\documentclass[../../../include/open-logic-section]{subfiles}

\begin{document}

\olfileid{nml}{com}{lin}

\olsection{Lemma Lindenbaum}

Lemma Lindenbaum netepake yen saben himpunan !!{formula}s
$\Sigma$-konsisten kalebu ing paling ora siji himpunan \emph{jangkep}
$\Sigma$-konsisten. Pambangunan model kanonik bakal nuduhake yen kanggo
saben himpunan $\Sigma$-konsisten jangkep~$\Delta$, ana sawijining
jagad ing model kanonik kang kabeh lan mung !!{formula}s
ing~$\Delta$ bener. Dadi Lemma Lindenbaum njamin yen saben himpunan
$\Sigma$-konsisten bener ing sawijining jagad ing model kanonik.

\begin{thm}[Lemma Lindenbaum]\ollabel{thm:lindenbaum}
  Yen $\Gamma$ $\Sigma$-konsisten, mula ana himpunan $\Sigma$-konsisten
  jangkep $\Delta$ kang ngluwihi~$\Gamma$.
\end{thm}

\begin{proof}
Upamane $!A_0$, $!A_1$, \dots{} iku dhaptar kang nyakup kabeh formula
ing basa iki (pangulangan diidinake). Upamane, wiwitana kanthi
ndhaptar $\Obj p_0$, banjur ing saben tahap $n \ge 1$ dhaptarna
formula-formula kang cacahe winates lan dawane~$n$ kanthi mung nggunakake
variabel ing antarane $\Obj p_0$, \dots,~$\Obj p_n$. Kita netepake
himpunan !!{formula}s $\Delta_n$ kanthi induksi marang~$n$, banjur
netepake $\Delta = \bigcup_n \Delta_n$. Kawitan, tetepa
$\Delta_0 = \Gamma$. Yen $\Delta_n$ wis ditepakake, kita netepake
$\Delta_{n+1}$ kaya mangkene:
\[
\Delta_{n+1} =
\begin{cases}
  \Delta_n \cup \{!A_n\}, & \text{yen $\Delta_n \cup \{ !A_n\}$
    iku $\Sigma$-konsisten;} \\
  \Delta_n \cup \{ \lnot !A_n\}, & \text{yen ora.}
\end{cases}
\]
Saiki tetepa $\Delta = \bigcup_{n=0}^\infty \Delta_n$.

Kita kudu nuduhake yen definisi iki pancen ngasilake himpunan
$\Delta$ kanthi sipat-sipat kang dibutuhake, yaiku $\Gamma \subseteq \Delta$
lan $\Delta$ iku $\Sigma$-konsisten jangkep.

Cetha yen $\Gamma \subseteq \Delta$, amarga $\Delta_0 \subseteq
\Delta$ miturut pambangunan lan $\Delta_0 = \Gamma$. Malah, $\Delta_n
\subseteq \Delta$ kanggo kabeh~$n$, amarga $\Delta$ iku gabungan
kabeh~$\Delta_n$. (Amarga ing saben langkah pambangunan kita nambah
!!a{formula} marang himpunan kang wis dibangun, $\Delta_n \subseteq
\Delta_{n+1}$; mula amarga $\subseteq$ transitif, $\Delta_n \subseteq
\Delta_{m}$ yen $n \le m$.) Ing saben tahap pambangunan, kita nambah
$!A_n$ utawa $\lnot !A_n$, lan saben !!{formula} katon (paling ora
sapisan) ing dhaptar kabeh~$!A_n$. Mula kanggo saben $!A$, $!A
\in \Delta$ utawa $\lnot !A \in \Delta$; kanthi definisi, $\Delta$
jangkep.

Pungkasan, kita kudu nuduhake yen $\Delta$ iku
$\Sigma$-konsisten. Kanggo iku, kita nuduhake yen (a) manawa $\Delta$
ora $\Sigma$-konsisten, sawijining $\Delta_n$ uga ora
$\Sigma$-konsisten, lan (b) kabeh $\Delta_n$ iku $\Sigma$-konsisten.

Upamane $\Delta$ ora $\Sigma$-konsisten. Banjur $\Delta
\Proves[\Sigma] \lfalse$, yaiku ana $!A_1$, \dots,~$!A_k \in
\Delta$ saengga $\Sigma \Proves !A_1 \lif (!A_2 \lif \cdots (!A_k
\lif \lfalse)\dots)$. Amarga $\Delta = \bigcup_{n=0}^\infty \Delta_n$,
saben $!A_i \in \Delta_{n_i}$ kanggo sawenehing~$n_i$. Pilihen $n$
minangka kang paling gedhe ing antarane indeks-indeks iki. Amarga
$n_i \le n$, $\Delta_{n_i} \subseteq \Delta_n$. Dadi kabeh
$!A_i$ ana ing sawijining~$\Delta_n$. Iki ateges $\Delta_n
\Proves[\Sigma] \lfalse$, yaiku $\Delta_n$ ora $\Sigma$-konsisten.

Kanggo nuduhake yen saben $\Delta_n$ $\Sigma$-konsisten,
kita nggunakake induksi prasaja marang~$n$. $\Delta_0 = \Gamma$,
lan kita wis nganggep $\Gamma$ $\Sigma$-konsisten. Dadi pratelan iki
bener kanggo $n = 0$. Saiki upamane pratelan iki bener kanggo $n$,
yaiku $\Delta_n$ $\Sigma$-konsisten. $\Delta_{n+1}$ iku
$\Delta_n \cup \{!A_n\}$ yen himpunan iki $\Sigma$-konsisten;
yen ora, iku $\Delta_n \cup \{\lnot!A_n\}$. Ing
kasus kapisan, $\Delta_{n+1}$ cetha $\Sigma$-konsisten. Nanging miturut
\olref[prf][con]{prop:consistencyfacts}\olref[prf][con]{prop:consistencyfacts-c},
$\Delta_n \cup \{!A_n\}$ utawa $\Delta_n \cup \{\lnot!A_n\}$
konsisten, mula $\Delta_{n+1}$ uga konsisten ing kasus liyane.
\end{proof}

\begin{cor}\ollabel{cor:provability-characterization}
  $\Gamma \Proves[\Sigma] !A$ yen lan mung yen $!A \in \Delta$ kanggo
  saben himpunan $\Sigma$-konsisten jangkep $\Delta$ kang ngluwihi $\Gamma$
  (kalebu nalika $\Gamma = \emptyset$, kang menehi karakterisasi liyane
  tumrap sistem modal $\Sigma$.)
\end{cor}

\begin{proof}
  Upamane $\Gamma \Proves[\Sigma] !A$, lan $\Delta$ sembarang himpunan
  $\Sigma$-konsisten jangkep kang ngluwihi $\Gamma$. Yen $!A
  \notin \Delta$, mula miturut maksimalitas $\lnot!A \in \Delta$, saengga
  $\Delta \Proves[\Sigma] !A$ (miturut monotonisitas) lan $\Delta
  \Proves[\Sigma] \lnot!A$ (miturut refleksivitas); mula $\Delta$
  ora konsisten. Kosok baline, yen $\Gamma \Proves/[\Sigma] !A$, mula
  $\Gamma \cup \{ \lnot!A\}$ $\Sigma$-konsisten, lan miturut
  Lemma Lindenbaum ana himpunan konsisten jangkep $\Delta$
  kang ngluwihi $\Gamma \cup \{ \lnot!A \}$. Amarga konsistensi, $!A
  \notin \Delta$.
\end{proof}

\end{document}
